\chapter{Evaluation}
\label{chapter:evaluation}

Following the design and implementation of the anonymization service prototype, this chapter provides a systematic evaluation of the system against the functional and nonfunctional requirements established in Chapter \ref{chapter:requirements}.  Ultimately, this evaluation determines the extent to which the service provides a legally and technically sound basis for measuring developer productivity in a research or production environment.

\section{Functional Requirements}

\subsubsection{FLR-1: Data Anonymization}

Evaluating the quality of an anonymization process requires a multifaceted analysis that balances technical robustness, legal compliance, and data utility. The "goodness" of the resulting dataset is primarily measured through formal anonymity models — such as k-anonymity — which provide mathematical metrics for the degree of protection against identity and attribute disclosure. Beyond these syntactic criteria, the effectiveness of the transformation must be validated against the three fundamental risk scenarios: singling out (isolating an individual), linkability (connecting records across datasets), and inference (deducing sensitive values from other attributes). A robust evaluation further involves the implementation of an attacker model, which simulates realistic adversaries with varying levels of background knowledge and resources to quantify the remaining residual risk. From a legal perspective, the anonymization is considered successful if it achieves factual anonymity, meaning that re-identification is either impossible or requires a disproportionate effort in terms of time, cost, and manpower. Finally, the quality of the process is intrinsically linked to data utility; a high-quality anonymization must minimize information loss to ensure that the dataset remains suitable for the intended analytical purpose while maintaining individual privacy.
\autocite[73--75,81--82,86--88,103--106,120-121]{Krebs2022}
\autocite[61--63]]{gmdsbvd2024praxishilfe}
\autocite[3--6,28--29]{schwartmann2022praxisleitfaden}

A comprehensive quantitative evaluation of every potential re-identification scenario or the implementation of advanced privacy models such as l-diversity and t-closeness is beyond the scope of this thesis. Instead, the prototype focuses on establishing a robust baseline for privacy by first pseudonymizing direct identifiers, which is achieved by hashing them into deterministic \acp{UUID} within the identity management system. Building upon this layer of protection, k-anonymity is achieved through the aggregation of data points based on organizational hierarchy. This is specifically realized by grouping individuals by their enrollment, a mechanism that links contributors to specific organizations over defined time intervals. By organizing data into these equivalence classes, the service ensures that each record becomes indistinguishable from at least k−1 other individuals, thereby effectively mitigating the risk of singling out specific developers while preserving the utility of group-level productivity metrics.

\subsubsection{FSR-1: Preserve Data}

The evaluation of FSR-1: Preserve Data reveals an inherent technical and legal conflict with the requirement for data anonymization (FLR-1), as the absolute preservation of all source data is impossible when specific attributes must be transformed or removed to protect individual privacy. In the current prototype, this requirement is addressed through a tiered approach: direct identifiers are pseudonymized by hashing them into deterministic \acp{UUID}, which remain resolvable through the SortingHat identity management system to maintain longitudinal traceability. To further protect privacy, individual contributors are generalized into equivalence classes based on their organizational enrollment, effectively achieving k-anonymity. While the retention of pseudonymized data is permissible and necessary during the scientific research and development phase of MECOIS, it may be mandatory to irreversibly delete the mapping "keys" in the final product to ensure compliance with the GDPR’s anonymization standards. Consequently, the preservation of data utility remains dynamic; indirect identifiers, such as Git commit messages, may also be subject to future deletion if they pose a disproportionate re-identification risk.

\subsubsection{FSR-2: Independent Deployment}

The prototype successfully achieves architectural and operational decoupling from the primary MECOIS framework. By design, the anonymization service operates as an isolated system that communicates via standard HTTP interfaces, allowing it to be deployed as a standalone component or as part of a fully integrated on-premise installation. From a security perspective, enclosing sensitive data within these independent systems establishes a strict isolation boundary, thereby reducing the risk of unauthorized data exposure and preventing sensitive payloads from leaking into neighboring host components or the primary analytics engine. This technical independence is primarily realized through the use of Docker Compose, which encapsulates the service and its specific environment within isolated containers, thereby preventing dependency conflicts and ensuring consistent behavior across different deployment environments. While the service functions as a proxy for the SortingHat, this decoupled architecture provides the necessary flexibility to scale the anonymization layer independently of the main analytics engine to meet varying performance demands.

\subsubsection{FSR-3: Deterministic vs. Non-Deterministic Pseudonym	Generation}

The prototype successfully prioritizes a deterministic approach to ensure the consistency of longitudinal data analysis. By utilizing the native \ac{UUID} generation mechanism within the SortingHat identity management system — which hashes direct identifiers such as names, emails, and usernames — the service ensures that the same individual is assigned a consistent pseudonym across different pipeline runs. This determinism is essential for the MECOIS framework's goal of linking cross-platform metrics from disparate sources like GitHub and Jira without revealing the developers' actual identities. However, the implementation is currently constrained by the use of the SHA-1 hashing algorithm, which presents cryptographic vulnerabilities such as a lack of resistance to brute-force attacks and an absence of salting. While this simplified method is functionally sufficient for the prototype's current research phase and supports the necessary traceability, the evaluation highlights that a more secure, salted hashing scheme will be required for the final product to more effectively mitigate re-identification risks.

\subsubsection{FIR-1: Seamless Integration}

The anonymization service successfully integrates into the MECOIS framework by replacing legacy processing steps with calls to the new server. Specifically, the previous, insufficient approach to anonymization by pseudonymizing direct identifiers in the Bronze-to-Silver pipeline has been replaced by invoking the \texttt{change\_identities} endpoint, ensuring that the data is anonymized as an integral part of the data flow. Similarly, the Silver-to-Silver pipeline has been adapted to utilize the \texttt{update\_identities} endpoint to maintain data consistency when changes occur within the identity management system. Furthermore, existing communication with the SortingHat system has been successfully redirected through the new service, which now acts as a proxy for requests such as \texttt{show-profile} or \texttt{find-profile}. This transition fulfills the requirement for minimal disruption to existing features while providing the system with fine-grained control over the personal information exposed to the analytics engine.

\section{Nonfunctional Requirements}

\subsubsection{NFSR-1: Containerized Service}

The requirement is fully met through the implementation of \texttt{Docker Compose}, which manages the service as part of a multi-container application via a unified YAML configuration. This containerized approach ensures environmental consistency and application isolation, effectively preventing dependency conflicts that could arise during the transition from development to production. Furthermore, the encapsulation of the application logic within a container provides a secure, isolated instance that utilizes its own namespaces and local file system to remain distinct from the host machine and other integrated services. By leveraging these containerization technologies, the prototype achieves a high-level abstraction that allows for the simultaneous initiation of all necessary microservices with a single command, fulfilling the demand for a portable and robust deployment model.

\subsubsection{NFSR-2: Horizontal Scalability and Statelessness}

The anonymization service is designed as a stateless component, fulfilling the architectural requirement for efficient horizontal scaling. By functioning primarily as a transformation engine for incoming DataFrames and acting as a proxy for the SortingHat API, the service avoids maintaining local session state between individual HTTP requests. This stateless architecture ensures that multiple instances of the service can be seamlessly deployed behind a load balancer to distribute the computational load during heavy pipeline jobs without risk of data inconsistency. Furthermore, the implementation of containerization via \texttt{Docker Compose} facilitates the rapid initiation of additional isolated instances, providing the necessary elasticity to handle varying processing demands within the MECOIS framework. This decoupled design allows the anonymization layer to scale its resources independently of the main analytics engine, thereby ensuring both performance and high availability.

\subsubsection{NFIR-1: Data Integrity}

The anonymization service effectively maintains the consistency of data structures during its integration with the broader MECOIS framework. This requirement is successfully addressed by the \texttt{utils.converter} module, which manages the serialization and deserialization of \texttt{PySpark DataFrames} and their corresponding schemas into \texttt{JSON} format for HTTP transmission. By preserving schema integrity, the system ensures that no data loss occurs during the communication between the decentralized system components, thereby maintaining the quality of the activity-based artifacts required for statistical analysis. This technical measure is critical for the reliable derivation of the Productivity Index, as it ensures that the high-dimensional metadata remains structurally sound and analytically useful throughout the entire anonymization pipeline.

\subsubsection{NFIR-2: Programming Language}

The requirement is fully satisfied, as the anonymization service is implemented entirely in \texttt{Python3}. This alignment with the existing MECOIS backend ensures seamless interoperability between the system components, particularly for the processing of \texttt{PySpark DataFrames} and the integration of the service into the multi-stage data pipeline. Furthermore, the use of \texttt{Python} allows the prototype to leverage standardized libraries, such as \texttt{http.server} for the backend implementation and \texttt{hashlib} for generating deterministic \acp{UUID}, while adhering to the project's mandated coding and license policies. By fulfilling this requirement, the service maintains technical consistency across the research framework, facilitating easier maintenance and future development.

\subsubsection{NFIR-3: Licence Policy}

The anonymization service adheres to the project's mandate to avoid the use of packages with copyleft or proprietary licenses. The prototype is implemented using \texttt{Python } and relies primarily on its standard libraries, such as \texttt{http.server} for the backend, which conform to permissive licensing standards. Furthermore, while the service interfaces with the SortingHat identity management system—which is governed by a \texttt{GPL-3.0} license — it maintains a strictly decoupled relationship by communicating exclusively via \texttt{HTTP}. This architectural separation ensures that the service's internal codebase remains independent of external copyleft requirements, thereby fulfilling the integration constraints established by the MECOIS contributing guidelines.

\subsubsection{NFIR-4: Styleguide and Clean Code}

The implementation of the anonymization service strictly adheres to the \texttt{MECOIS} contributing guidelines to ensure high software quality and long-term maintainability. The codebase is developed in alignment with the \texttt{PEP8} standard. By fulfilling these requirements, the service ensures technical consistency with the rest of the \texttt{MECOIS} backend, facilitating a seamless review and merging process into the main branch.
